\subsubsection{请求归属、信用与地址保护}
四项在途记录保存已发命令的读写归属及结束标志；命令握手增加归属队列，响应握手按序移除。读返回先进入 16 字缓冲，消费者接受后才减少缓冲占用。于是读空间判断同时考虑已返回但尚未消费的数据和已发出但尚未返回的读请求，不能只看当前 FIFO 内已有的数据量。

设 $O$ 为在途命令数、$R$ 为读缓冲占用、$R_p$ 为未返回读数，则发出新请求需要满足对应容量约束：在途槽未满，或本周期恰有响应释放一槽；读请求还要求 $R+R_p<16$，或本周期恰有读数据被消费。写请求从 8 字写缓冲取数，在读写均可发送时优先写出。相同周期的入队和出队按净变化更新计数，因此连续事务不必为每次槽位释放额外插入空拍。

\begin{figure}[H]\centering
\begin{tikzpicture}
\node[box,text width=31mm] (r) at (0,1.1) {读请求资格\\已提交边界、读空间};
\node[box,text width=31mm] (w) at (0,-.7) {写请求资格\\写缓冲中已有数据};
\node[box,text width=30mm] (s) at (4.2,.2) {在途信用检查\\写优先选择};
\node[box,text width=30mm] (h) at (8.5,.2) {保持寄存器\\ICB 命令接口};
\node[box,text width=45mm] (o) at (6.5,-2) {4 项归属记录\\响应返回后按读写归属退休};
\draw[arr](r)--(s);\draw[arr](w)--(s);\draw[arr](s)--(h);
\draw[arr](h.east)--++(.5,0)|-(o.east);
\draw[arr,accent](o.west)--(3,-2)--(3,-1)--(s.south);
\end{tikzpicture}
\caption{DMA 的请求选择和响应归属。返回归属记录用于流量控制，不增加共享 SRAM 的访问端口。}
\end{figure}

描述符要求基址按 4 字节对齐、起始位置不小于 \code{0xa0000000}，传输字数在 1 到 131072 之间，结束地址不超过 \code{0xa0080000}。结束地址以扩展位宽计算，避免 32 位加法回绕绕过范围检查。输入可读边界还受实时提交字数限制，因此一个描述符合法并不表示其全部内容已由装载端写入。

\subsubsection{反压期间的稳定性}
当 cmd\_valid 为 1、cmd\_ready 为 0，命令已经对外呈现但尚未被接受，地址、数据、方向及结束属性必须保持。保持寄存器将这些值锁定，直到握手；此时即使内部选择器出现新的读写请求，也不能替换已呈现的命令。发生 abort 后同样先处理已经呈现的命令和在途响应，再回到空闲状态。

\begin{figure}[H]\centering
\begin{tikzpicture}[x=1.65cm,y=.70cm]
\foreach \i in {0,...,5}{\draw[gray!30](\i,-.2)--(\i,-4.5);\node[note]at(\i+.5,.35){边沿 \i};}
\node[anchor=east,note]at(-.2,-.65){cmd\_valid};
\node[anchor=east,note]at(-.2,-1.65){cmd\_ready};
\node[anchor=east,note]at(-.2,-2.65){地址、数据};
\node[anchor=east,note]at(-.2,-3.65){命令接收};
\draw[ink,thick](0,-.9)--(1,-.9)--(1,-.4)--(5,-.4)--(5,-.9)--(6,-.9);
\draw[ink,thick](0,-1.4)--(1,-1.4)--(1,-1.9)--(4,-1.9)--(4,-1.4)--(6,-1.4);
\draw[fill=accent!15,draw=ink](1,-2.95)rectangle(5,-2.35);\node[note]at(3,-2.65){同一命令，保持稳定};
\draw[ink,thick](0,-3.9)--(4,-3.9)--(4,-3.4)--(5,-3.4)--(5,-3.9)--(6,-3.9);
\node[note]at(3,-4.65){ready 恢复后完成一次握手；周期仅用于说明协议};
\end{tikzpicture}
\caption{ICB 反压握手示意，非仿真波形或实测时延。高电平重叠处接受命令。}
\end{figure}

已记录的两种错误注入分别在下游就绪和不就绪时触发非法输出索引。两次均在 5 周期内完成排空，命令与响应总数保持相等。这个结果支持所覆盖注入位置的排空行为，不意味着所有错误位置、任意总线停顿长度都固定为 5 周期；若下游长期不返回响应，清空时间仍由真实响应进度决定。
